% ==========================================
\section{Dettaglio dello Stack Tecnologico}
% ==========================================

Il progetto tiene client e server ben separati, ma coerenti tra loro: i due lati condividono i tipi TypeScript e le stesse regole di validazione.

\subsection{Client (Frontend) -- \texttt{client/package.json}}
\begin{table}[H]
\centering
\arrayrulecolor{SlateGray}
\renewcommand{\arraystretch}{1.3}
\begin{tabularx}{\textwidth}{|l|c|X|}
\hline
\rowcolor[HTML]{3B3766}
\textcolor{white}{\textbf{Libreria}} & \textcolor{white}{\textbf{Versione}} & \textcolor{white}{\textbf{Scopo}} \\ \hline
\textbf{\texttt{@angular/core} (\dots)} & \textasciicircum{}21.2 & Framework SPA (architettura Standalone + Signals) \\ \hline
\textbf{\texttt{@ionic/angular}} & \textasciicircum{}8.0 & Componenti UI ibridi e approccio mobile-first \\ \hline
\textbf{\texttt{ionicons}} & \textasciicircum{}7.0 & Set di icone vettoriali standardizzate \\ \hline
\textbf{\texttt{@capacitor/core}} & 8.x & Shell nativa per la build e il deploy su dispositivi mobili \\ \hline
\textbf{\texttt{echarts}} & \textasciicircum{}6.1 & Motore di rendering grafico (candele, linee, aree, volumi) \\ \hline
\textbf{\texttt{@jsverse/transloco}} & \textasciicircum{}7.5 & Gestione dell'internazionalizzazione dinamica (IT/EN) \\ \hline
\textbf{\texttt{rxjs}} & \textasciitilde{}7.8 & Gestione asincrona e programmazione reattiva \\ \hline
\textbf{\texttt{zod}} & \textasciicircum{}3.25 & Validazione rigorosa dei form e dei payload lato client \\ \hline
\textbf{\texttt{zone.js} / \texttt{tslib}} & \textasciitilde{}0.15 / \textasciicircum{}2.8 & Runtime Angular e helper TypeScript \\ \hline
\end{tabularx}
\end{table}

\textbf{Tooling e Dipendenze di Sviluppo (devDependencies):}
\begin{itemize}[itemsep=-2pt]
    \item \textbf{Build/CLI:} \texttt{@angular/build}, \texttt{@angular/cli}, \texttt{@angular/\allowbreak{}compiler-cli}, \texttt{@angular/\allowbreak{}language-service} (\textasciicircum{}21), \texttt{@ionic/\allowbreak{}angular-toolkit}, \texttt{@capacitor/cli}.
    \item \textbf{Linting:} \texttt{eslint} (\textasciicircum{}9), i pacchetti \texttt{@angular-eslint/*} e \texttt{@typescript-eslint/*}, oltre a plugin specifici per import, JSDoc e arrow functions.
    \item \textbf{Testing:} \texttt{karma}, \texttt{jasmine-core}, e relativi reporter/launcher per la suite di test.
    \item \textbf{Linguaggio:} \texttt{typescript} \textasciitilde{}5.9.
\end{itemize}

\subsection{Server (Backend) -- \texttt{server/package.json}}
\begin{table}[H]
\centering
\arrayrulecolor{SlateGray}
\renewcommand{\arraystretch}{1.3}
\begin{tabularx}{\textwidth}{|l|c|X|}
\hline
\rowcolor[HTML]{3B3766}
\textcolor{white}{\textbf{Libreria}} & \textcolor{white}{\textbf{Versione}} & \textcolor{white}{\textbf{Scopo}} \\ \hline
\textbf{\texttt{express}} & \textasciicircum{}4.19 & Web framework e gestore dell'infrastruttura API REST \\ \hline
\textbf{\texttt{pg}} & \textasciicircum{}8.12 & Driver nativo per il client PostgreSQL (connessione a Supabase) \\ \hline
\textbf{\texttt{jsonwebtoken}} & \textasciicircum{}9.0 & Generazione e verifica dei token per l'autenticazione stateless \\ \hline
\textbf{\texttt{bcryptjs}} & \textasciicircum{}2.4 & Funzioni di hashing crittografico per la protezione delle password \\ \hline
\textbf{\texttt{zod}} & \textasciicircum{}3.23 & Validazione rigorosa degli input e dei payload in ingresso \\ \hline
\textbf{\texttt{nodemailer}} & \textasciicircum{}9.0 & Modulo di invio email transazionali tramite protocollo SMTP \\ \hline
\textbf{\texttt{cors}} & \textasciicircum{}2.8 & Configurazione e gestione delle policy CORS \\ \hline
\textbf{\texttt{dotenv}} & \textasciicircum{}16.4 & Iniezione e gestione sicura delle variabili d'ambiente \\ \hline
\end{tabularx}
\end{table}

\textbf{Tooling e Dipendenze di Sviluppo (devDependencies):}
\begin{itemize}[itemsep=-2pt]
    \item \textbf{Esecuzione:} \texttt{tsx} (\textasciicircum{}4.16) per l'esecuzione nativa e il \textit{watch} dei file TypeScript in ambiente di sviluppo.
    \item \textbf{Linguaggio e Tipizzazione:} \texttt{typescript} (\textasciicircum{}5.5) e pacchetti \texttt{@types/} per Node, Express, pg, JWT, bcryptjs, cors e nodemailer.
\end{itemize}

\subsection{Modulo condiviso}
Il progetto ha una cartella \texttt{shared/} che fa da ponte tra frontend e backend: non contiene codice eseguibile, solo modelli e interfacce TypeScript usati da entrambi (importabili con l'alias \texttt{\$shared}). \textbf{Zod} è l'unica libreria usata su tutti e due i lati, così le regole di validazione sono scritte una volta sola. Il server usa Express~4 (\textasciicircum{}4.19), ancora il più compatibile con i middleware disponibili.

% ==========================================
\section{Attori del Sistema}
% ==========================================
Gli attori sono i profili, umani o automatici, che interagiscono con il sistema; ognuno ha permessi diversi.

\begin{table}[H]
\centering
\arrayrulecolor{SlateGray}
\renewcommand{\arraystretch}{1.3}
\begin{tabularx}{\textwidth}{|l|c|X|X|}
\hline
\rowcolor[HTML]{3B3766}
\textcolor{white}{\textbf{Attore}} & \textcolor{white}{\textbf{Tipo}} & \textcolor{white}{\textbf{Descrizione}} & \textcolor{white}{\textbf{Ruolo e Permessi nel Sistema}} \\ \hline
\textbf{Utente Standard (Trader)} & Umano & Chi usa la piattaforma per analizzare e simulare. & Accede con una sola sessione attiva. Gestisce i propri workspace, mette e collega i widget, compra e vende nel Broker (con commissione), cura il profilo e partecipa alla classifica (condivisione e recensioni). \\ \hline
\textbf{Amministratore (Admin)} & Umano & Utente con privilegi in più, che supervisiona la piattaforma. & Ha i permessi dell'utente standard più il Pannello Admin: aggiunge o toglie titoli dal listino, elimina account, forza l'aggiornamento dei prezzi e vede i ricavi dell'app. \\ \hline
\textbf{Scheduler (Sistema)} & Automatico & Processo che gira in background sul server Node.js. & Svolge compiti a tempo senza intervento umano: aggiorna prezzi e storico da Yahoo Finance, e gestisce la scadenza dei token di reset. \\ \hline
\textbf{Provider Dati} & Servizio esterno & API di terze parti usata dal backend. & Fornisce i dati reali (quotazioni e serie storiche) per i widget azionari e i grafici. \\ \hline
\end{tabularx}
\end{table}

% ==========================================
\section{Modello Dati e Struttura del Database}
% ==========================================

I dati di DeQ sono salvati in un database PostgreSQL ospitato su Supabase. Lo schema tiene separati lo stato finanziario dell'utente, i workspace salvati e lo storico delle transazioni.

\begin{figure}[H]
    \centering
    \includegraphics[width=\textwidth]{Immagini/drawsql.png}
    \caption{Diagramma Entità-Relazione (ER) del database di sistema generato tramite DrawSQL.}
    \label{fig:er-diagram}
\end{figure}

\subsection{Entità principali}
Sono le tabelle al centro dell'applicazione, quelle che reggono i dati dell'utente e del mercato.

\begin{table}[H]
\centering
\arrayrulecolor{SlateGray}
\renewcommand{\arraystretch}{1.3}
\begin{tabularx}{\textwidth}{|l|X|p{3.5cm}|X|}
\hline
\rowcolor[HTML]{3B3766}
\textcolor{white}{\textbf{Entità}} & \textcolor{white}{\textbf{Descrizione}} & \textcolor{white}{\textbf{Attributi Critici}} & \textcolor{white}{\textbf{Implicazioni Logiche / Relazioni}} \\ \hline
\textbf{User} & L'account registrato, con i suoi dati e il suo stato economico. & \texttt{id}, \texttt{email}, \texttt{credit}, \texttt{is\_\allowbreak{}premium}, \texttt{token\_\allowbreak{}version} & Tabella centrale. \texttt{token\_\allowbreak{}version} impedisce i login simultanei. È in relazione 1:N con workspace, transazioni e posizioni. \\ \hline
\textbf{Asset} & Il titolo (azione quotata o personalizzata) che si può seguire e comprare. & \texttt{ticker}, \texttt{name}, \texttt{last\_\allowbreak{}price}, \texttt{currency} & Identificato dal ticker. Alimenta i grafici ed è ciò che si compra nel broker. \\ \hline
\textbf{Holding} & Una posizione: quante quote di un titolo ha un utente. & \texttt{id}, \texttt{user\_\allowbreak{}id}, \texttt{ticker}, \texttt{quantity}, \texttt{avg\_\allowbreak{}price} & Collega User e Asset (relazione M:N). Salva il prezzo medio d'acquisto, usato per il calcolo del guadagno. \\ \hline
\textbf{Transaction} & Il registro (solo aggiunte) di ogni movimento di saldo o stato. & \texttt{id}, \texttt{user\_\allowbreak{}id}, \texttt{type}, \texttt{amount}, \texttt{created\_\allowbreak{}at} & Segna il motivo (acquisto, vendita, ricarica, Premium), utile per lo storico e per l'admin. \\ \hline
\textbf{Workspace} & Una scheda della canvas, con la disposizione dei nodi. & \texttt{id}, \texttt{user\_\allowbreak{}id}, \texttt{name}, \texttt{camera\_\allowbreak{}zoom}, \texttt{widget} (\texttt{JSONB}) & Salva la posizione della vista e l'intera rete di widget e collegamenti in un campo JSON, per il salvataggio automatico. \\ \hline
\end{tabularx}
\end{table}

\subsection{Entità di supporto}
Tabelle di servizio: classifica, cache dei prezzi, sicurezza e contabilità.

\begin{table}[H]
\centering
\arrayrulecolor{SlateGray}
\renewcommand{\arraystretch}{1.3}
\begin{tabularx}{\textwidth}{|l|X|X|X|}
\hline
\rowcolor[HTML]{3B3766}
\textcolor{white}{\textbf{Entità}} & \textcolor{white}{\textbf{Descrizione}} & \textcolor{white}{\textbf{Funzione Architetturale}} & \textcolor{white}{\textbf{Allineamento e Vincoli}} \\ \hline
\textbf{Leaderboard Entry} & La scheda che un utente rende pubblica per la classifica. & Ordina la classifica in base al guadagno o alla perdita delle posizioni (prezzo attuale meno prezzo medio d'acquisto). & Relazione N:1 con User. Il vincolo unico (\texttt{user\_\allowbreak{}id}, \texttt{label}) permette a un utente di partecipare con più schede. \\ \hline
\textbf{Leaderboard Review} & La recensione che un utente lascia sulla scheda di un altro. & È la parte sociale: commenti e confronto tra utenti. & Due chiavi esterne in \texttt{CASCADE}: verso l'utente recensito (\texttt{entry\_\allowbreak{}user\_\allowbreak{}id}) e verso l'autore (\texttt{author\_\allowbreak{}id}). \\ \hline
\textbf{Market Cache} & Copia temporanea delle risposte di Yahoo Finance. & Riduce le chiamate esterne ed evita di superarne i limiti. & Il dato vale circa 20 minuti (in base a \texttt{last\_\allowbreak{}fetched}). Indicizzata per \texttt{ticker} e \texttt{timeframe}. \\ \hline
\textbf{Price History} & Storico dei prezzi registrati nel tempo dal sistema. & Serve al nastro di mercato e al broker per calcolare le variazioni. & Tabella a sé (solo aggiunte). Si aggancia ai titoli solo per ticker, così lo storico resta anche se un titolo viene tolto. \\ \hline
\textbf{Password Reset} & Il record temporaneo per il recupero password. & Conserva l'hash del token e la scadenza (1 ora) per validare il link via email. & Relazione N:1 con User, con un flag \texttt{used} per non riutilizzare lo stesso token. \\ \hline
\textbf{App Revenue} & Registro dei ricavi (fittizi) dell'app. & Somma gli incassi da Premium (€10) e dalle commissioni (1\%) per i grafici dell'admin. & Tabella isolata (nessuna chiave esterna), così i dati contabili restano anche se un utente viene eliminato. \\ \hline
\end{tabularx}
\end{table}

\subsection{Relazioni e vincoli}
La coerenza dei dati è affidata direttamente ai vincoli del database:

\begin{itemize}
    \item \textbf{Cancellazione a cascata (\texttt{ON DELETE CASCADE}):} eliminando un utente si cancellano anche i suoi dati collegati: posizioni (\texttt{holdings}), transazioni, workspace e token di recupero password.
    \item \textbf{Riferimento azzerato (\texttt{ON DELETE SET NULL}):} le colonne che indicano chi ha aggiunto un titolo (\texttt{added\_by} in \texttt{assets}) o chi ha fatto un'azione (\texttt{actor\_id} in \texttt{asset\_events}) vengono messe a NULL se quell'admin viene eliminato, così lo storico resta.
    \item \textbf{Vincoli di unicità (\texttt{UNIQUE}):} la coppia \texttt{(user\_id, ticker)} in \texttt{holdings} evita di spezzare una stessa azione in più righe; la coppia \texttt{(user\_id, label)} in \texttt{leaderboard\_\allowbreak{}entries} permette a un utente di avere più schede in classifica, ma senza duplicati.
\end{itemize}

\subsection{Tabelle indipendenti}
Alcune tabelle sono volutamente slegate dalle chiavi esterne, così restano valide anche se il listino cambia:

\begin{itemize}
    \item \textbf{Storico e cache dei prezzi:} \texttt{market\_\allowbreak{}cache} e \texttt{price\_\allowbreak{}history} si agganciano ai titoli solo tramite il \texttt{ticker} (testo). Se un'azione viene tolta dal listino, lo storico e la cache restano comunque disponibili.
    \item \textbf{Contabilità:} \texttt{app\_\allowbreak{}revenue} (ricavi da commissioni e Premium) è una tabella a sé: i conti dell'app non cambiano anche se l'utente che li ha generati viene eliminato.
\end{itemize}